\section{Determinant suppliers controlled by the branch exponent}
\label{sec:branch-exponent}

\paragraph{Provenance and logical status.}
The Abel-fibre strategy originates in Schoen's published construction
\cite[\S2, pp.~12--17]{G34:Schoen}.  The arbitrary-branching version was
located in the unrefereed OpenAI manuscript \cite[Lemma~3.4]{Root:Gale},
and the improvement from the character order to the common inertia
exponent was located in \cite[\S2]{Root:Rstation66}.  These are credited
as leads, not assumed as theorems.  The proof below reconstructs the
geometric argument, including the direction of the returning
correspondence, from the stated published inputs.  In particular it
does not use the global CM Hodge conjecture, 032, or a Severn theorem.
The quantitative degree of this construction is a separate assertion;
no numerical detector bound follows just from the word ``supplied.''

Here a sub-Hodge structure of the cohomology of a smooth projective
variety is \emph{supplied by} a smooth projective variety $F$ if it is
contained in the sum of the images of finitely many algebraic
correspondences from powers of $F$, with Tate twists.  Coefficients may
first be extended to a number field.  The rational correspondence span
is stable under scalar conjugation, tensor products, algebraic
projectors, and composition.  These assertions follow directly by
external product and composition of correspondences.

\begin{theorem}[HD-FS1: the inertia-exponent cyclic supplier]
\label{HD-FS1}
Let $q:C\to B$ be a connected cyclic cover of smooth projective
complex curves, of degree $m\ge2$, and let $\chi$ be a faithful character
of its deck group.  Suppose every inertia order divides the positive
integer $M$.  Then
\[
 D_\chi=\det H^1(C,\mathbf Q(\zeta_m))_\chi
\]
is supplied by the Jacobian of the degree-$M$ Fermat curve.  Unramified
covers need only point suppliers.  The assertion includes a
zero-dimensional eigenspace, whose determinant is the unit.
\end{theorem}
\begin{proof}
Write $b=g(B)$ and let $D=t_1+\cdots+t_r$ be the reduced branch divisor.
The relation in the punctured surface group, followed by its cyclic
quotient, shows that $r\ne1$.  Put
\[
              v=\dim H^1(C,\mathbf Q(\zeta_m))_\chi=2b-2+r.
\]
Indeed the character part of $q_*\mathbf Q(\zeta_m)$ is a rank-one
local system on $B\setminus D$, extended by zero at $D$; its Euler
characteristic is $2-2b-r$.  Its degrees zero and two vanish because
the cover is connected and the character is nontrivial.  If $v=0$
there is nothing to prove.  Assume $v>0$.

\emph{The alternating line as quotient cohomology.}
Let $\Delta$ be the cyclic deck group, and let $N$ be the kernel of
the product homomorphism
\[
          \Delta^v\rtimes\mathfrak S_v\longrightarrow\Delta.
\]
The map is a homomorphism since $\Delta$ is abelian.  Set
$X=C^v$ and $Z=X/N$; there is a finite map
$\pi:Z\to\operatorname{Sym}^v B$ with residual group $\Delta$.
Finite quotient pullback identifies rational cohomology with the
invariants: proper base change kills higher direct images, and
averaging makes invariants exact.  In particular $H^v(Z)$ is pure,
since it injects into $H^v(X)$.
The $\chi$-eigenspace of $H^v(Z)$ maps to $D_\chi$ in $H^v(X)$.
To verify this, each curve factor must have degree one and character
$\chi$; the geometric permutation of degree-one classes contributes
its Koszul sign, so invariance under permutations is ordinary
alternation.  The resulting exterior power is a line.

On $U=\operatorname{Sym}^v(B\setminus D)$, the residual cover is a
$\Delta$-torsor, including at divisors with repeated points.  In a
small disk, sheet labels modulo $N$ remember only their product,
and elementary symmetric functions give coordinates on the symmetric
power.  At a divisor containing a branch point the character stalk
is zero, since its stabilizer contains nontrivial inertia.  Thus
$(\pi_*\mathbf Q(\zeta_m))_\chi$ is extension by zero from a
rank-one local system on $U$.  Its monodromy on a moving divisor
is the product of the original cover monodromies along the moving
point paths.  This description is valid across diagonals by the
local torsor calculation.

\emph{Localization at one Abel fibre.}
First suppose $r>0$, so $v\ge2b$.  Let
\[
 a:\operatorname{Sym}^v B\longrightarrow\operatorname{Pic}^v B,
 \qquad L_0=K_B(D),\qquad
 T=(\pi^{-1}(|L_0|))_{\rm red}.
\]
For $L\ne L_0$ the evaluation map
$H^0(B,L)\to\bigoplus_i L|_{t_i}$ is surjective.  Its obstruction
is dual to the section space of the nontrivial degree-zero bundle
$K_B\otimes L^{-1}(D)$.  Its kernel has dimension $b-1$.
Consequently
\[
                |L|\cap U\simeq\mathbf A^{b-1}\times
                                      (\mathbf C^*)^{r-1}.
\]
The rank-one local system on this product has a nontrivial circle
monodromy.  A transverse loop about an evaluation hyperplane winds
the product of the nearby zero coordinates once around the relevant
branch point, so its monodromy is that inertia character.  This also
handles multiple zeros, by the product of roots in a local monic
equation.  A nontrivial rank-one local system on a circle has cellular
complex $k\xrightarrow{\alpha-1}k$ with $\alpha\ne1$; tensoring
these complexes proves vanishing on the torus, and duality gives
compact-support vanishing.  Proper base change and Leray therefore
give the restriction isomorphism
\[
                    H^v(Z)_\chi\xrightarrow{\sim}H^v(T)_\chi.
\tag{FS.1}
\]
If $b=0$, the Picard variety is a point and the same assertion is
immediate.

The residue sequence identifies
\[
 H^0(B,K_B(D))\simeq W\oplus R,
 \quad W=H^0(B,K_B),\quad
 R=\{(w_1,\ldots,w_r):\textstyle\sum_iw_i=0\},
\]
after choosing a vector-space splitting.  The evaluation hyperplanes
are $w_i=0$ in $\mathbf P(W\oplus R)$.  For $r\ge3$ they are
distinct.  Over its function field the residual torsor has a Kummer
equation
\[
                    z^m=c h^m\prod_{i=2}^r(w_i/w_1)^{a_i}.
\tag{FS.2}
\]
Here $c\in\mathbf C^*$ and $h$ is a rational function.  To justify
the assertion globally, its valuations away from the hyperplanes
are divisible by $m$.  Choose the integers $a_i$ to match the
hyperplane valuations modulo $m$; the degree-zero condition matches
the remaining hyperplane.  The residual divisor is $mE$ and the
torsion-free Picard group of projective space makes $E$ principal.
The local transverse-loop calculation above shows that the order
of $\exp(2\pi i a_i/m)$ divides $M$.  Equivalently,
\[
                             m\mid Ma_i.
\tag{FS.3}
\]
This is the step that replaces the character order by the inertia
exponent.

Let $F_M^{r-2}\subset\mathbf P^{r-1}$ be the Fermat hypersurface
$\sum u_i^M=0$ and use its finite coordinate-power map
$s:F_M^{r-2}\to\mathbf P(R)$, $[u_i]\mapsto[u_i^M]$.
The smooth projective bundle
\[
        Y=\mathbf P_{F_M^{r-2}}
              ((W\otimes\mathcal O)\oplus s^*\mathcal O(-1))
\tag{FS.4}
\]
maps generically finitely onto $\mathbf P(W\oplus R)$.
After this pullback each factor in (FS.2) is an $m$th power by
(FS.3).  The generic torsor splits.  Rational lifts to its components
extend to morphisms: a rational lift to a finite scheme over the
base is regular on a normal source by integral closure of every
local ring.  Taking finitely many copies of components of $Y$ gives
a proper surjection $Y'\to T$.  Every component of $T$ dominates the
base: its translates cover every fibre, which is one $\Delta$-orbit.
Thus this surjection misses no component.  For $r=2$ the evaluation
hyperplanes coincide, the Kummer expression has no nonconstant
factor, and $T$ is instead dominated by a finite union of projective
spaces.  This is also the dimension-zero Fermat case.

Shioda--Katsura's published Theorem~1.7 and Proposition~2.4
\cite{Root:SK} give algebraic correspondence supply of the full
cohomology of $F_M^d$ by powers of $J(F_M^1)$, for every $d$.
Indeed their degree-$M$ rational map from a product of lower
Fermat hypersurfaces is resolved by their displayed smooth blowup;
the blowup formula, induction, and push-pull by a generically finite
map give surjectivity on cohomology.  Projective-bundle cohomology
is a direct sum of Tate twists of the base via powers of its
tautological class.  Hence the entire cohomology of $Y'$ has the
required supply.

For completeness, if $r=0$, then $v=2b-2$ and $b\ge2$.  The
nontrivial cover character defines a rank-one local system $E$ on
$\operatorname{Pic}^vB$, and its pullback by $a$ is the character
local system on the symmetric power.  The Abel fibres are
$\mathbf P^{b-2}$ except over $K_B$, where the fibre is
$\mathbf P^{b-1}$.  The powers of a relatively ample class show
that the even direct images up to degree $2b-4$ are constant;
the remaining direct image is a rank-one skyscraper in degree
$2b-2$ at $K_B$.  Tensor with $E$ and use the same torus
acyclicity on the Picard variety.  This proves (FS.1), now with
$T$ a split finite cover of the simply connected $|K_B|$.
It is a union of projective spaces and is supplied by points.

\emph{Returning from detection to supply.}
Deligne's published Proposition~8.2.5, p.~39, together with
Theorem~8.2.4(iii) \cite{Root:DeligneIII}, gives
\[
 \ker\bigl(H^v(T,\mathbf Q)\longrightarrow H^v(Y',\mathbf Q)\bigr)
                       =W_{v-1}H^v(T,\mathbf Q).
\tag{FS.5}
\]
The pure line from $H^v(Z)$ has nonzero restriction by (FS.1)
and cannot meet this kernel, by strictness of mixed Hodge structures.
Choose a component $Y_0$, of dimension $d$, on which a nonzero
$\alpha\in H^v(Z)_\chi$ survives.  Poincare duality gives
$\beta\in H^{2d-v}(Y_0)$ with
$\int_{Y_0}p^*\alpha\,\beta\ne0$, where $p:Y_0\to Z$.
In the branched case $2d-v=r-2\ge0$; in the unramified case it is
zero.  The full-cohomology supply just proved applies to $\beta$.

Take a component of $Y_0\times_Z X$ dominating $Y_0$ and resolve
it projectively to $\widetilde Y$, with maps
$h:\widetilde Y\to Y_0$ and $l:\widetilde Y\to X$.
Its cycle $(h,l)_*[\widetilde Y]$ defines the correspondence
$l_*h^*$ and
\[
       \int_X Q^*\alpha\, l_*h^*\beta
           =\deg(h)\int_{Y_0}p^*\alpha\,\beta\ne0,
\tag{FS.6}
\]
where $Q:X\to Z$.  No supply of the resolution's cohomology is
required.  Apply the inverse-character projector for
$\Delta^v\rtimes\mathfrak S_v$, a finite linear combination of
algebraic automorphism graphs.  In degree $v$ its image is precisely
$D_{\chi^{-1}}$, by the earlier Koszul calculation, and the
nonzero pairing (FS.6) is unchanged.  Thus that line is supplied.
Conjugating coefficient scalars supplies $D_\chi$ as well.
All maps are algebraic correspondences, and the only nontrivial
fixed source is $J(F_M^1)$.
\end{proof}

\begin{theorem}[HD-FS2: finite-group transfer with one Fermat exponent]
\label{HD-FS2}
Let a finite group $G$ act on a smooth projective curve, allowing a
disjoint union of connected curves permuted transitively.  Suppose
every inertia order in its quotient cover divides $M$.  Over a
splitting field, the determinant of every irreducible multiplicity
space in $H^1$ is supplied by $J(F_M^1)$.  If all powers of this
Jacobian have the rational Hodge conjecture, every balanced rational
sum of these determinant lines is algebraic.
\end{theorem}
\begin{proof}
First take a connected cover.  Integral Brauer induction
\cite[\S10.5, Theorems~19--20]{G34:Serre} gives an identity
\[
 V\oplus R_-\simeq R_+,
 \qquad R_\pm=\bigoplus_j
        (\operatorname{Ind}_{H_j}^G\chi_j)^{\oplus n_{j,\pm}},
 \quad n_{j,\pm}\ge0,
\tag{FS.7}
\]
with linear characters $\chi_j$.  Frobenius reciprocity and quotient
pullback identify their multiplicity spaces with the corresponding
faithful character spaces for
$C/\ker\chi_j\to C/H_j$.  The inertia generators of each of these
intermediate covers are powers of conjugates of the original
generators.  Their orders therefore divide $M$.  HD-FS1 applies.
The trivial character contributes the determinant of $H^1(C/H_j)$,
which is a Tate line supplied by a top polarization class.

Apply $\operatorname{Hom}_G(-,H^1(C))$ to (FS.7) and take
determinants.  It gives
$D_+=D_V\otimes D_-$.  Tensor with the determinant for the
conjugate of the minus representation and contract the two minus
determinants by the nondegenerate polarization pairing.  This
contraction is algebraic: the $H^1$-Kunneth projector and the
diagonal of a pointed curve realize the pairing and its inverse,
and the finite group algebra and matrix units realize the required
summands.  Consequently the tensor of the supplied plus and
conjugate-minus lines maps nontrivially onto $D_V$, with a Tate
shift.  This cancels an actual dual factor; it does not take a root
of a cycle class.

For a disconnected $G$-cover, choose one connected component,
with stabilizer $H$.  Its full cohomology induces to that of the
disjoint union.  Restrict $V$ to $H$, decompose it into irreducibles,
and apply the connected assertion and determinant tensor products.
Again the inertia orders divide $M$.

Finally a rational Hodge class in an algebraic image of a sum of
polarizable pure Hodge structures has a rational Hodge preimage:
take the orthogonal complement of the kernel with respect to a
polarization.  Under the stated Hodge-conjecture hypothesis this
preimage is algebraic on a power of the fixed Jacobian.  Its
algebraic-correspondence image is algebraic.  Extending coefficients
and then descending the rational cycle-class image deals with the
individual character lines.  No assertion is made that an
unbalanced character determinant is itself a Hodge class.
\end{proof}

\paragraph{The $E_8$ application and its remaining numerical step.}
Every branch order in HD-E8-02 and D-E8T-01 divides $12$.
Thus the fixed source in HD-FS2 is $J(F_{12})$, irrespective of
the much larger exponent of $\mu_4W(E_8)$ or of the orders of
Brauer's inducing characters.  A controlled elliptic decomposition
of this fixed Jacobian and explicit bounds for the correspondence
degrees will be recorded separately.  The present theorem proves
the supply statement, not a numerical value of $\delta_W$.
